\documentclass{situationist-preprint}

\title{Meaning Without Projects}
\msid{EXI-01-c}
\msfamily{EXI-01}
\msdiscipline{Meaning, Death and the Absurd}
\mstype{Alternative proof}
\mskeywords{meaning in life; projects; objective worth; engagement; availability}

\begin{document}
\maketitle

\begin{abstract}
EXI-01-a establishes that a life of positive meaning is available to every agent capable of intentional action. Its proof rests on the assumption that every such agent is able to engage in at least one project of objective worth (Assumption~3). That assumption is open to doubt for agents whose circumstances deny them the opportunities that projects ordinarily require. We give a proof of the same conclusion that does not assume it. We begin not from projects but from objects of engagement, such as persons and the world as an object of understanding, show that such objects are within the reach of every agent, and show that sustained engagement with one of them constitutes a project in the sense relevant to the definition of meaning. Assumption~3 is thereby derived rather than assumed.
\end{abstract}

\section{The premise to be avoided}

EXI-01-a defines a life as meaningful to the degree that it consists in active engagement with projects of objective worth (Definition~1), assumes that there are such projects (Assumption~2), and assumes that every agent capable of intentional action is able to engage in at least one of them (Assumption~3). Theorem~4 follows at once.

Assumption~3 carries the weight of the universality claim, and it is the premise a critic will press. Projects, as ordinarily understood, are extended undertakings with aims: a career, a cause, a work to be made. Many lives offer little room for them. Prisoners, those who labor without respite, the very ill and the very poor may have no project of objective worth on which they can embark, and \citet{Wolf2010} acknowledges that circumstances can leave people with few opportunities for meaning. If Assumption~3 is read in this ordinary sense, it is doubtful. The proof below does not use it. It keeps Definition~1 and Assumption~2 of EXI-01-a.

\section{The alternative route}

\begin{definition}\label{d:object}
An \emph{object of engagement} is anything at which an agent can direct intentional activity and about which she can care: a person, a living creature, a practice, an inquiry, or a part of the world as an object of understanding. It is \emph{worthwhile} if it has or serves value that does not depend on the agent's attitudes toward it (Definition~2.2 of EXI-17-a).
\end{definition}

The list follows the objects that accounts of meaning have most often recognized: persons, the pursuit of understanding and the appreciation of value \citep{Metz2013}.

\begin{assumption}[Reach]\label{a:reach}
Every agent capable of intentional action can direct some intentional activity, with care, at a worthwhile object of engagement: another person, or some part of the world she can come to understand.
\end{assumption}

\begin{lemma}[Projects from engagement]\label{l:proj}
If an agent directs intentional activity with care at a worthwhile object of engagement, and sustains it, her engagement is a project of objective worth in the sense of Definition~1 of EXI-01-a.
\end{lemma}

\begin{proof}[Argument]
What Definition~1 requires of a project is that it be something the agent engages in, actively, and that its value not depend on the attitudes of those engaged in it (Assumption~2 of EXI-01-a). Sustained, caring activity directed at an object has the first feature: it is an undertaking with a direction, extended over time, that the agent can pursue well or badly. It has the second when its object is worthwhile, for the value of caring for a person or of coming to understand something derives from the value of the person or of the understanding, not from the agent's attitude. Nothing in Definition~1 requires that a project have a completion condition, a public product or a social form; \citet{Wolf2010} counts relationships of love and the pursuit of understanding among her central cases.
\end{proof}

\begin{theorem}\label{t:main}
For every agent capable of intentional action, a life of positive meaning is available. The proof does not use Assumption~3 of EXI-01-a.
\end{theorem}

\begin{proof}
By Assumption~\ref{a:reach}, the agent can direct intentional activity with care at a worthwhile object of engagement. By Lemma~\ref{l:proj}, sustaining that activity is engagement in a project of objective worth. A life that includes such engagement consists to a positive extent in active engagement with a project of objective worth, and by Definition~1 of EXI-01-a it is meaningful to a positive degree.
\end{proof}

Assumption~3 of EXI-01-a now follows as a corollary of Assumption~\ref{a:reach} and Lemma~\ref{l:proj}, and is no longer a premise.

\section{Why the route is better supported}

Assumption~3 asserts the availability of projects; Assumption~\ref{a:reach} asserts the availability of objects. The difference matters because opportunities for projects in the ordinary sense depend on resources, institutions and freedom, which are unequally distributed, while worthwhile objects are present in every human situation. Almost every agent lives among other persons, whose good has value independent of her attitudes, and every agent lives in a world some part of which she can attend to and understand. The prisoner can care for a fellow prisoner and can come to understand something; the exhausted laborer can attend to her children; the bedridden patient can attend to the people who visit her and to what she reads or hears. EXI-17-a reaches a parallel result for achievement: a life can be meaningful through engagement without success (Theorem~3.3 of EXI-17-a), and through appreciation without any product (Lemma~3.2 of EXI-17-a). The present argument applies the same point to opportunity: meaning does not wait on the availability of a project, because engagement makes one.

\section{Objections}

\begin{objection}
Lemma~\ref{l:proj} trivializes the notion of a project. If attending to a visitor is a project, the word has lost its content.
\end{objection}

\begin{reply}
The lemma requires sustained activity, care and a worthwhile object. A visit attended to once is not a project; a relationship of care sustained over time is. What the lemma drops are features of projects that are irrelevant to meaning on Definition~1: a completion condition, a public product and institutional form. Their irrelevance is what EXI-17-a shows for achievement.
\end{reply}

\begin{objection}
Assumption~\ref{a:reach} is Assumption~3 in other words.
\end{objection}

\begin{reply}
It entails Assumption~3 given Lemma~\ref{l:proj}, as it must if the conclusion is to be reached. But the two premises differ in their support. Assumption~3, read as a claim about projects in the ordinary sense, is open to counterexample from deprived circumstances; Assumption~\ref{a:reach} is supported by the ubiquity of persons and of the world as an object of understanding, and the cases that threatened Assumption~3 satisfy it.
\end{reply}

\begin{objection}
Some agents capable of intentional action are wholly isolated and wholly unable to learn anything further.
\end{objection}

\begin{reply}
An agent capable of intentional action acts on something, and anything she acts on is part of a world she can attend to. The case of an agent for whom nothing worthwhile is within reach of attention is, if possible at all, a case of the absence of agency rather than of its exercise, and falls outside the scope of both proofs.
\end{reply}

\section{Conclusion}

The universal availability of meaning does not depend on the universal availability of projects. It depends on something more widely distributed: the presence, in every human life, of worthwhile things to care about and to understand. Engagement with them is a project of objective worth, and a life that contains it has meaning. This reaches the conclusion of EXI-01-a without its third premise.

\begin{thebibliography}{9}\small
\bibitem[Metz(2013)]{Metz2013} Metz, T. (2013). \emph{Meaning in Life: An Analytic Study}. Oxford University Press.
\bibitem[Wolf(2010)]{Wolf2010} Wolf, S. (2010). \emph{Meaning in Life and Why It Matters}. Princeton University Press.
\end{thebibliography}

\end{document}
